% ============================================================
\chapter{Probability Spaces}
\label{ch:spaces}
% ============================================================

The calculus of probabilities was born in the gambling houses of the
seventeenth century. Blaise Pascal and Pierre de Fermat, in their famous
1654 correspondence, sought to solve the ``problem of points'': how to
fairly divide the stakes of a game that is interrupted. But it took nearly
three centuries for probability theory to acquire rigorous mathematical
foundations. It was Andrei Kolmogorov who, in 1933 in his
\emph{Grundbegriffe der Wahrscheinlichkeitsrechnung}, laid down the
definitive axioms: a probability is nothing but a measure of total mass $1$
on a measurable space.

This first chapter builds these foundations stone by stone: the sample space
$\Omega$ of all possibilities, the $\sigma$-algebra $\mathcal{F}$ of
observable events, and the probability measure $\PP$ that assigns each event
a degree of likelihood.

\begin{intuition}
This opening chapter lays the mathematical foundations of probability theory.
We define the notion of a probability space $(\Omega, \mathcal{F}, \PP)$,
present the Kolmogorov axioms, and establish the first properties of the
probability measure.
\end{intuition}

% ------------------------------------------------------------
\section{Random Experiments and Sample Spaces}
% ------------------------------------------------------------

Before any formalization, let us ask the naive question: what is a random
experiment? It is an experiment whose outcome we cannot predict, but whose
set of possible outcomes we can describe. A die roll, tomorrow's stock
price, the lifetime of a transistor --- all situations where uncertainty is
irreducible, yet the structure of that uncertainty is accessible.

\begin{definition}[Random experiment]
A \textbf{random experiment} is an experiment whose outcome cannot be predicted
with certainty before it is performed, but for which the set of all possible
outcomes can be described.
\end{definition}

\begin{definition}[Sample space]
The \textbf{sample space} is the set $\Omega$ of all possible outcomes of a
random experiment. Each element $\omega \in \Omega$ is called a
\textbf{sample point} or \textbf{outcome}.
\end{definition}

\begin{example}
\begin{enumerate}
    \item \textbf{Roll of a die:} $\Omega = \{1, 2, 3, 4, 5, 6\}$.
    \item \textbf{Coin toss:} $\Omega = \{\text{Heads}, \text{Tails}\}$.
    \item \textbf{Lifetime of a component:} $\Omega = [0, +\infty)$.
    \item \textbf{Infinite sequence of coin tosses:}
          $\Omega = \{0, 1\}^{\N} = \{(\omega_1, \omega_2, \ldots) : \omega_i \in \{0,1\}\}$.
\end{enumerate}
\end{example}

\begin{remark}
The sample space $\Omega$ may be finite, countably infinite, or uncountable.
The choice of $\Omega$ depends on the model; it is not unique for a given experiment.
\end{remark}

% ------------------------------------------------------------
\section{$\sigma$-Algebras and Events}
% ------------------------------------------------------------

When $\Omega$ is finite or countable we can usually take every subset to be an
event. For uncountable spaces such as $\R$, we must restrict to a well-behaved
class of ``measurable'' sets.

\begin{definition}[$\sigma$-algebra]
\label{def:sigma_algebra}
A \textbf{$\sigma$-algebra} (or $\sigma$-field) on $\Omega$ is a collection
$\mathcal{F} \subseteq \mathcal{P}(\Omega)$ satisfying:
\begin{enumerate}[label=(\roman*)]
    \item $\Omega \in \mathcal{F}$;
    \item if $A \in \mathcal{F}$, then $A^c \in \mathcal{F}$
          \quad (closure under complementation);
    \item if $(A_n)_{n \geq 1}$ is a sequence in $\mathcal{F}$, then
          $\displaystyle\bigcup_{n=1}^{\infty} A_n \in \mathcal{F}$
          \quad (closure under countable unions).
\end{enumerate}
\end{definition}

\begin{proposition}[Immediate properties]
Let $\mathcal{F}$ be a $\sigma$-algebra on $\Omega$. Then:
\begin{enumerate}[label=(\alph*)]
    \item $\emptyset \in \mathcal{F}$;
    \item $\mathcal{F}$ is closed under countable intersections:
          if $A_n \in \mathcal{F}$ for all $n$, then
          $\bigcap_{n=1}^{\infty} A_n \in \mathcal{F}$;
    \item $\mathcal{F}$ is closed under set differences: if $A, B \in \mathcal{F}$,
          then $A \setminus B \in \mathcal{F}$.
\end{enumerate}
\end{proposition}

\begin{proof}
(a) $\emptyset = \Omega^c \in \mathcal{F}$ by (i) and (ii).
(b) By De Morgan's laws, $\bigcap A_n = \bigl(\bigcup A_n^c\bigr)^c$,
which belongs to $\mathcal{F}$ by (ii) and (iii).
(c) $A \setminus B = A \cap B^c \in \mathcal{F}$ by (ii) and (b).
\end{proof}

\begin{definition}[Event]
An \textbf{event} is an element of the $\sigma$-algebra $\mathcal{F}$, i.e.\ a
subset of $\Omega$ to which a probability can be assigned. We say event $A$
\textbf{occurs} if the outcome $\omega$ belongs to $A$.
\end{definition}

\begin{example}[Fundamental $\sigma$-algebras]
\begin{enumerate}
    \item \textbf{Trivial $\sigma$-algebra:} $\mathcal{F} = \{\emptyset, \Omega\}$.
          The smallest $\sigma$-algebra on $\Omega$.
    \item \textbf{Discrete $\sigma$-algebra:} $\mathcal{F} = \mathcal{P}(\Omega)$.
          The largest $\sigma$-algebra on $\Omega$.
    \item \textbf{Generated $\sigma$-algebra:} if $\mathcal{C}$ is a collection of
          subsets of $\Omega$, then $\sigma(\mathcal{C})$ denotes the smallest
          $\sigma$-algebra containing $\mathcal{C}$.
\end{enumerate}
\end{example}

\begin{definition}[Borel $\sigma$-algebra]
The \textbf{Borel $\sigma$-algebra} on $\R$, denoted $\mathcal{B}(\R)$, is the
$\sigma$-algebra generated by the open subsets of $\R$:
\[
    \mathcal{B}(\R) = \sigma\bigl(\{O \subseteq \R : O \text{ is open}\}\bigr).
\]
One can show that $\mathcal{B}(\R) = \sigma\bigl(\{(-\infty, a] : a \in \R\}\bigr)
= \sigma\bigl(\{(a, b) : a < b\}\bigr)$.
\end{definition}

% ------------------------------------------------------------
\section{The Kolmogorov Axioms}
% ------------------------------------------------------------

\begin{definition}[Probability measure]
\label{def:prob_measure}
Let $(\Omega, \mathcal{F})$ be a measurable space. A \textbf{probability measure}
is a function $\PP : \mathcal{F} \to [0,1]$ satisfying the \textbf{Kolmogorov axioms}:
\begin{enumerate}[label=\bfseries (K\arabic*)]
    \item \textbf{Non-negativity:} for every $A \in \mathcal{F}$,
          $\PP(A) \geq 0$.
    \item \textbf{Normalisation:} $\PP(\Omega) = 1$.
    \item \textbf{$\sigma$-additivity:} for every sequence $(A_n)_{n \geq 1}$
          of \emph{mutually disjoint} events (i.e.\ $A_i \cap A_j = \emptyset$
          for $i \neq j$):
          \[
              \PP\!\left(\bigcup_{n=1}^{\infty} A_n\right)
              = \sum_{n=1}^{\infty} \PP(A_n).
          \]
\end{enumerate}
\end{definition}

\begin{definition}[Probability space]
The triple $(\Omega, \mathcal{F}, \PP)$ is called a \textbf{probability space}.
\end{definition}

\begin{warning}[title=\textbf{Finite vs.\ $\sigma$-additivity}]
Finite additivity ($\PP(A \cup B) = \PP(A) + \PP(B)$ when $A \cap B = \emptyset$)
follows from (K3). The $\sigma$-additivity axiom is strictly stronger: it guarantees
consistency when passing to limits and is essential for the limit theorems of
Chapters~9--11.
\end{warning}

% ------------------------------------------------------------
\section{Properties of the Probability Measure}
% ------------------------------------------------------------

\begin{theorem}[Fundamental properties]
\label{thm:prob_properties}
Let $(\Omega, \mathcal{F}, \PP)$ be a probability space. For all
$A, B \in \mathcal{F}$:
\begin{enumerate}[label=(\roman*)]
    \item $\PP(\emptyset) = 0$.
    \item $\PP(A^c) = 1 - \PP(A)$.
    \item If $A \subseteq B$ then $\PP(A) \leq \PP(B)$ \quad (monotonicity).
    \item $\PP(A \cup B) = \PP(A) + \PP(B) - \PP(A \cap B)$.
    \item $\PP(A) \leq 1$.
\end{enumerate}
\end{theorem}

\begin{proof}
\begin{enumerate}[label=(\roman*)]
    \item Set $A_1 = \Omega$ and $A_n = \emptyset$ for $n \geq 2$. The sequence is
    disjoint and $\bigcup A_n = \Omega$. By (K3):
    $1 = \PP(\Omega) = \PP(\Omega) + \sum_{n \geq 2}\PP(\emptyset)$,
    so $\PP(\emptyset) = 0$.

    \item $\Omega = A \cup A^c$ (disjoint), so $1 = \PP(A) + \PP(A^c)$.

    \item $B = A \cup (B \setminus A)$ (disjoint), so
    $\PP(B) = \PP(A) + \PP(B \setminus A) \geq \PP(A)$.

    \item $A \cup B = A \cup (B \setminus A)$ and
    $B = (A \cap B) \cup (B \setminus A)$, giving
    $\PP(B \setminus A) = \PP(B) - \PP(A \cap B)$, and the result follows.

    \item From (ii): $\PP(A) = 1 - \PP(A^c) \leq 1$.
\end{enumerate}
\end{proof}

\begin{theorem}[Inclusion--exclusion formula]
\label{thm:inclusion_exclusion}
For events $A_1, \ldots, A_n \in \mathcal{F}$:
\[
    \PP\!\left(\bigcup_{i=1}^{n} A_i\right)
    = \sum_{i} \PP(A_i)
    - \sum_{i < j} \PP(A_i \cap A_j)
    + \sum_{i < j < k} \PP(A_i \cap A_j \cap A_k)
    - \cdots
    + (-1)^{n+1}\, \PP(A_1 \cap \cdots \cap A_n).
\]
\end{theorem}

\begin{corollary}[Boole's inequality / sub-additivity]
For any sequence $(A_n)_{n \geq 1}$ of events:
\[
    \PP\!\left(\bigcup_{n=1}^{\infty} A_n\right)
    \leq \sum_{n=1}^{\infty} \PP(A_n).
\]
\end{corollary}

% ------------------------------------------------------------
\section{Continuity of the Probability Measure}
% ------------------------------------------------------------

\begin{theorem}[Monotone continuity]
\label{thm:continuity}
Let $(\Omega, \mathcal{F}, \PP)$ be a probability space.
\begin{enumerate}[label=(\roman*)]
    \item \textbf{Continuity from below:} if $A_1 \subseteq A_2 \subseteq \cdots$
    (increasing), then
    \[
        \PP\!\left(\bigcup_{n=1}^{\infty} A_n\right)
        = \lim_{n \to \infty} \PP(A_n).
    \]
    \item \textbf{Continuity from above:} if $A_1 \supseteq A_2 \supseteq \cdots$
    (decreasing), then
    \[
        \PP\!\left(\bigcap_{n=1}^{\infty} A_n\right)
        = \lim_{n \to \infty} \PP(A_n).
    \]
\end{enumerate}
\end{theorem}

\begin{proof}
(i) Define $B_1 = A_1$ and $B_n = A_n \setminus A_{n-1}$ for $n \geq 2$. The $B_n$
are mutually disjoint and $\bigcup_{n=1}^{N} B_n = A_N$. By $\sigma$-additivity:
\[
    \PP\!\left(\bigcup_{n=1}^{\infty} A_n\right)
    = \PP\!\left(\bigcup_{n=1}^{\infty} B_n\right)
    = \sum_{n=1}^{\infty} \PP(B_n)
    = \lim_{N \to \infty} \sum_{n=1}^{N} \PP(B_n)
    = \lim_{N \to \infty} \PP(A_N).
\]

(ii) Apply (i) to the increasing sequence $A_1^c \subseteq A_2^c \subseteq \cdots$
and use $\PP(A_n^c) = 1 - \PP(A_n)$.
\end{proof}

% ------------------------------------------------------------
\section{Constructing Probability Measures}
% ------------------------------------------------------------

\subsection{Uniform probability on a finite set}

\begin{definition}[Uniform probability]
If $\Omega$ is a finite set, the \textbf{uniform probability} (or equally likely
model) is defined by:
\[
    \PP(A) = \frac{\abs{A}}{\abs{\Omega}}
    = \frac{\text{number of favourable outcomes}}{\text{total number of outcomes}}.
\]
\end{definition}

\begin{intuition}[title=\textbf{When is the uniform model appropriate?}]
The assumption of equally likely outcomes is only legitimate when the symmetry of
the problem justifies it (fair die, random draw, etc.). Never invoke it without
verification!
\end{intuition}

\begin{example}[The birthday problem]
In a group of $n$ people, what is the probability that at least two share the
same birthday? Assume 365 equally likely days.

The sample space is $\Omega = \{1, \ldots, 365\}^n$ with $\abs{\Omega} = 365^n$.
The complementary event ``all birthdays are distinct'' has cardinality
$365 \times 364 \times \cdots \times (365 - n + 1)$. Hence:
\[
    \PP(\text{at least one match})
    = 1 - \frac{365!}{(365-n)!\, 365^n}.
\]
For $n = 23$ we obtain $\PP \approx 0.507$: more likely than not!
\end{example}

\subsection{Probability defined by a discrete mass function}

\begin{proposition}
If $\Omega$ is countable and $(p_\omega)_{\omega \in \Omega}$ is a family of reals
with $p_\omega \geq 0$ for all $\omega$ and $\sum_{\omega \in \Omega} p_\omega = 1$,
then
\[
    \PP(A) = \sum_{\omega \in A} p_\omega, \quad A \subseteq \Omega,
\]
defines a probability measure on $(\Omega, \mathcal{P}(\Omega))$.
\end{proposition}

\subsection{Construction on $\R$: continuous densities}

\begin{proposition}
Let $f : \R \to [0, +\infty)$ be an integrable function with
$\int_{-\infty}^{+\infty} f(x)\, dx = 1$. Then
\[
    \PP(A) = \int_A f(x)\, dx, \quad A \in \mathcal{B}(\R),
\]
defines a probability measure on $(\R, \mathcal{B}(\R))$. The function $f$ is
called a \textbf{probability density function} (pdf).
\end{proposition}

% ------------------------------------------------------------
\section{The First Borel--Cantelli Lemma}
% ------------------------------------------------------------

\begin{theorem}[First Borel--Cantelli lemma]
\label{thm:borel_cantelli_1}
Let $(A_n)_{n \geq 1}$ be a sequence of events. If
$\sum_{n=1}^{\infty} \PP(A_n) < \infty$, then
\[
    \PP\!\left(\limsup_{n \to \infty} A_n\right) = 0,
\]
where $\displaystyle\limsup_{n \to \infty} A_n
= \bigcap_{n=1}^{\infty} \bigcup_{k=n}^{\infty} A_k$ is the event
``infinitely many $A_n$ occur''.
\end{theorem}

\begin{proof}
Set $B_n = \bigcup_{k=n}^{\infty} A_k$. The sequence $(B_n)$ is decreasing and
$\limsup A_n = \bigcap B_n$. By continuity from above:
\[
    \PP(\limsup A_n) = \lim_{n \to \infty} \PP(B_n)
    \leq \lim_{n \to \infty} \sum_{k=n}^{\infty} \PP(A_k) = 0,
\]
since the tail of a convergent series tends to zero.
\end{proof}

\begin{remark}
The \textbf{second} Borel--Cantelli lemma (a partial converse requiring independence)
will be stated in Chapter~2.
\end{remark}

% ------------------------------------------------------------
\section{Exercises}
% ------------------------------------------------------------

\begin{exercise}
Show that the intersection of two $\sigma$-algebras is a $\sigma$-algebra.
Is the union of two $\sigma$-algebras always a $\sigma$-algebra? Give a
counterexample.
\end{exercise}

\begin{exercise}
Let $\Omega = \{1,2,3,4,5,6\}$ (fair die). Write down the $\sigma$-algebra
generated by the partition $\bigl\{\{1,2\}, \{3,4\}, \{5,6\}\bigr\}$ and
count its elements.
\end{exercise}

\begin{exercise}
A fair die is rolled twice. Find the probability that:
\begin{enumerate}[label=(\alph*)]
    \item the sum equals 7;
    \item the sum is at least 10;
    \item both faces are the same.
\end{enumerate}
\end{exercise}

\begin{exercise}
Three cards are drawn successively without replacement from a standard deck of
52 cards. Find the probability of getting exactly two aces.
\end{exercise}

\begin{exercise}
Let $(A_n)$ be a decreasing sequence of events with $\PP(A_n) = 1/n$.
What is $\PP\bigl(\bigcap_{n=1}^{\infty} A_n\bigr)$? Justify your answer.
\end{exercise}

\begin{exercise}[Borel--Cantelli]
Let $(A_n)$ be a sequence of events with $\PP(A_n) = 1/n^2$.
Show that $\PP(\limsup A_n) = 0$. Now consider $\PP(A_n) = 1/n$: can you
still conclude?
\end{exercise}

\begin{exercise}
Prove the inclusion--exclusion formula (Theorem~\ref{thm:inclusion_exclusion})
by induction on $n$.
\end{exercise}

\begin{exercise}[Derangements]
Place $n$ letters into $n$ envelopes at random (one letter per envelope).
Let $D_n$ be the number of letters in the correct envelope. Using
inclusion--exclusion, show that
\[
    \PP(D_n = 0) = \sum_{k=0}^{n} \frac{(-1)^k}{k!}
    \xrightarrow[n \to \infty]{} e^{-1} \approx 0.368.
\]
\end{exercise}
